\documentclass[11pt]{article}

\usepackage[margin=1in]{geometry}
\usepackage{amsmath, amssymb, amsfonts}
\usepackage{booktabs}
\usepackage{longtable}
\usepackage{array}
\usepackage{hyperref}

\title{Code-Level Explainer for the iSFS Logistic Training Run}
\author{}
\date{\today}

\begin{document}
\maketitle

\section{Purpose of This Document}
This document explains the \emph{code-level} quantities that appear in the
training run and in the generated results tables. The earlier iSFS explainer was
mainly theoretical. This one answers a different question:
\begin{quote}
\emph{What do the variables and parameters in the code mean, where are they
used, and how do they relate to the block-coordinate-descent training loop?}
\end{quote}

The active implementation uses the iSFS structure from Section 3 of
\texttt{Multi\_source\_BWM.pdf}, replaces least squares with logistic/BCE loss,
and adds the trust-region cost update from \texttt{notes.tex}.

\section{Which Data Were Used in the Reported Run?}
The run
\begin{verbatim}
./run_training.sh --max-iter 10 --block-a-max-iter 40
\end{verbatim}
used the real \emph{blockwise-missing} files, not the complete files. In
\texttt{run\_training.py}, the real-data loader reads
\begin{align*}
X^{(1)} &: \texttt{data/source1\_clinical\_blockwise\_missing.csv}, \\
X^{(2)} &: \texttt{data/source2\_imaging\_blockwise\_missing.csv}, \\
X^{(3)} &: \texttt{data/source3\_proteomics\_blockwise\_missing.csv}, \\
y &: \texttt{data/labels.csv}.
\end{align*}

So the reported training metrics came from the blockwise-missing three-source
dataset with $500$ samples.

\section{The Main Training Objects}
For each source combination $m$, the model uses
\begin{align}
z_m = \sum_{i \in m} \alpha_m^i X_m^i \beta^i.
\end{align}
Here:
\begin{itemize}
\item $X_m^i$ is the source-$i$ feature block restricted to the samples in group
$m$,
\item $\beta^i$ is the source-level coefficient vector for source $i$,
\item $\alpha_m^i$ is the profile-specific source weight for source $i$ inside
group $m$.
\end{itemize}

For BCE from logits, the per-sample loss is
\begin{align}
\ell(z_{m,j}, y_{m,j}) = \log(1 + e^{z_{m,j}}) - y_{m,j} z_{m,j}.
\end{align}

With class costs $C$, the training group loss is
\begin{align}
f_m^C(\alpha,\beta,C)
=
\frac{1}{n_m}
\sum_{j \in G_m}
C_{y_{m,j}} \, \ell(z_{m,j}, y_{m,j}),
\end{align}
and the fixed-cost Block A objective is
\begin{align}
\Phi(\alpha,\beta;C)
=
\frac{1}{|pf|}
\sum_{m \in pf}
f_m^C(\alpha,\beta,C)
 \;+\; \lambda_\beta R_\beta(\beta).
\end{align}

In the active path,
\begin{align}
R_\beta(\beta) = \sum_i \|\beta^i\|_1,
\qquad
R_\alpha(\alpha_m) = \|\alpha_m\|_1 \le r_\alpha.
\end{align}

\section{Three Nested Iteration Levels}
The single most important source of confusion in the code is that there are
\textbf{three nested iteration levels}. They are not the same.

\subsection{Loop Schema}
\[
\boxed{
\begin{array}{l}
\textbf{Outer BCD loop } t = 1,\dots,T_{\mathrm{outer}} \\
\quad \text{Current accepted state: } (C^{(t)}, \alpha^{(t)}, \beta^{(t)}) \\
\quad \textbf{Block A solve with fixed } C^{(t)} \\
\qquad \textbf{Alpha subproblem loop } u = 1,\dots,T_\alpha \\
\qquad\qquad \beta \text{ fixed, update } \alpha \text{ by projected gradient} \\
\qquad \textbf{Beta subproblem loop } v = 1,\dots,T_\beta \\
\qquad\qquad \alpha \text{ fixed, update } \beta \text{ by proximal gradient} \\
\qquad \text{Repeat alpha/beta alternation up to } T_A \text{ times} \\
\quad \text{Compute class-loss vector } L(\alpha,\beta) \\
\quad \text{Solve trust-region cost step } d_t \\
\quad C_{\mathrm{trial}} = C^{(t)} + d_t \\
\quad \text{Re-solve Block A at } C_{\mathrm{trial}} \\
\quad \text{Compare current objective and trial objective} \\
\quad \text{Accept or reject the trial cost step using } \rho_t
\end{array}
}
\]

\subsection{What the Iteration Parameters Mean}
\begin{longtable}{>{\raggedright\arraybackslash}p{0.20\textwidth} >{\raggedright\arraybackslash}p{0.18\textwidth} >{\raggedright\arraybackslash}p{0.54\textwidth}}
\toprule
\textbf{Parameter} & \textbf{Loop Level} & \textbf{Meaning} \\
\midrule
\endhead
\texttt{max\_iter} & Outer BCD & Maximum number of outer cost-update rounds. One outer round means: fix $C$, solve Block A, propose a trial cost, re-solve Block A at the trial cost, then accept or reject the trial step. \\
\texttt{block\_a\_max\_iter} & Block A & Maximum number of alpha/beta alternations while $C$ is fixed. This is the loop that repeatedly calls the alpha solver and beta solver until the fixed-cost Block A objective stabilizes. \\
\texttt{alpha\_subproblem\_max\_iter} & Alpha subproblem & Maximum number of projected-gradient steps for solving one profile-specific alpha subproblem while beta is fixed. \\
\texttt{beta\_subproblem\_max\_iter} & Beta subproblem & Maximum number of proximal-gradient steps for solving the beta subproblem while alpha is fixed. \\
\texttt{tol} & Outer BCD & Outer stopping tolerance. It is used to terminate when the trust-region cost step is negligible or when the accepted outer improvement becomes tiny. \\
\texttt{block\_a\_tol} & Block A & Tolerance for deciding that the fixed-cost Block A objective has stabilized. \\
\texttt{alpha\_subproblem\_tol} & Alpha subproblem & Tolerance for deciding that the alpha iterate or the alpha objective has stabilized. \\
\texttt{beta\_subproblem\_tol} & Beta subproblem & Tolerance for deciding that the beta iterate or the beta objective has stabilized. \\
\bottomrule
\end{longtable}

\paragraph{Outer BCD vs. Block A.}
The difference is:
\begin{align*}
\text{Outer BCD} &:\quad \text{updates the cost vector } C, \\
\text{Block A} &:\quad \text{updates } (\alpha,\beta) \text{ while } C \text{ is fixed.}
\end{align*}
So \texttt{max\_iter} controls \emph{how many times the code tries to update
$C$}, whereas \texttt{block\_a\_max\_iter} controls \emph{how many
alpha/beta-alternation rounds are used to approximately solve the fixed-cost
subproblem}.

\section{Where the Learning Rates Enter}
The paper gives the optimization structure, but the code still needs a
numerical method to solve the alpha and beta subproblems. The current
implementation uses first-order iterative solvers, so it introduces the
step-size parameters
\begin{align}
\eta_\alpha = \texttt{learning\_rate\_alpha},
\qquad
\eta_\beta = \texttt{learning\_rate\_beta}.
\end{align}

These are \textbf{optimizer controls}, not new model parameters.

\subsection{Alpha Step}
For one profile $m$, the code computes a gradient
\begin{align}
\nabla_{\alpha_m} f_m^C(\alpha_m,\beta),
\end{align}
then takes a gradient step
\begin{align}
\alpha_m^{\mathrm{trial}}
=
\alpha_m - \eta_\alpha \nabla_{\alpha_m} f_m^C(\alpha_m,\beta),
\end{align}
and then projects back onto the $\ell_1$ ball:
\begin{align}
\alpha_m^{\mathrm{next}}
=
\Pi_{\|\cdot\|_1 \le r_\alpha}
\left(
\alpha_m^{\mathrm{trial}}
\right).
\end{align}

\subsection{Beta Step}
The beta step is analogous, except that the beta regularizer is treated through
a proximal map:
\begin{align}
\beta^{i,\mathrm{trial}}
&=
\beta^i - \eta_\beta \nabla_{\beta^i} g(\beta), \\
\beta^{i,\mathrm{next}}
&=
\mathrm{prox}_{\eta_\beta \lambda_\beta \|\cdot\|_1}
\left(
\beta^{i,\mathrm{trial}}
\right).
\end{align}

\paragraph{Interpretation.}
If the learning rate is too large, the iterative solver may overshoot. If it is
too small, the solver may move too slowly. So these are numerical controls for
the code implementation of the paper's subproblems.

\section{Two Different Balls: Do Not Mix Them Up}
There are \textbf{two different constraint geometries} in the code.

\subsection{Alpha \texorpdfstring{$\ell_1$}{l1} Ball}
The quantity
\begin{align}
r_\alpha = \texttt{alpha\_l1\_radius}
\end{align}
controls the constraint
\begin{align}
\|\alpha_m\|_1 \le r_\alpha
\qquad \text{for every profile } m.
\end{align}

This is \textbf{not} the trust-region ball for the cost vector. It belongs to
the alpha update only. In your run,
\begin{align}
r_\alpha = 1.0.
\end{align}
That is why the results table reported
\begin{verbatim}
Alpha Max L1: 1.000000
\end{verbatim}
many times: the projected alpha solution is living on the boundary of the alpha
$\ell_1$ ball.

\subsection{Cost Trust-Region Ball}
The cost vector $C$ is updated inside a \emph{different} Euclidean trust region.
That part uses:
\begin{align}
\Delta_t &= \text{working trust-region radius at iteration } t, \\
\Delta_{\max} &= \texttt{max\_trust\_region\_radius}, \\
C^0 &= \texttt{baseline\_cost\_vector}.
\end{align}

The local trust-region constraint is
\begin{align}
\frac{1}{2}\|d_t\|_2^2 \le \Delta_t,
\end{align}
while the global baseline-ball constraint is
\begin{align}
\frac{1}{2}\|C^{(t)} + d_t - C^0\|_2^2 \le \Delta_{\max}.
\end{align}

So:
\begin{align*}
\text{alpha ball} &:\ \ell_1 \text{ constraint on } \alpha_m, \\
\text{cost trust region} &:\ \ell_2 \text{ ball for updates of } C.
\end{align*}

\section{The Cost Vectors: Initial, Baseline, and Final}
The run report shows three cost-related objects.

\subsection{Initial Cost Vector}
The initial cost vector is the starting point
\begin{align}
C^{(0)} = \texttt{initial\_cost\_vector}.
\end{align}
In your run it was
\begin{align}
C^{(0)} = [0.5,\ 0.5].
\end{align}
The code uses this vector to solve the \emph{first} fixed-cost Block A problem.

\subsection{Baseline Cost Vector}
The baseline cost vector is
\begin{align}
C^0 = \texttt{baseline\_cost\_vector}.
\end{align}
This is the \emph{global reference center} of the large trust-region ball. It
does not move from one outer iteration to the next unless the user chooses a
different baseline at the start.

If the user does not supply a separate baseline, the code sets
\begin{align}
C^0 = C^{(0)}.
\end{align}
That is exactly what happened in your run, so the initial cost vector and the
baseline cost vector happened to be the same numerical vector:
\begin{align}
C^{(0)} = C^0 = [0.5,\ 0.5].
\end{align}

\paragraph{Important distinction.}
Even when the two vectors start equal, they play different roles:
\begin{align*}
C^{(0)} &: \text{starting current cost}, \\
C^0 &: \text{fixed baseline center for the global trust-region ball.}
\end{align*}

\paragraph{In your run.}
The command
\begin{verbatim}
./run_training.sh --max-iter 10 --block-a-max-iter 40
\end{verbatim}
used
\begin{align}
\texttt{initial\_cost\_vector} &= [0.5,\ 0.5], \\
\texttt{baseline\_cost\_vector} &= [0.5,\ 0.5].
\end{align}
So numerically, your run started with
\begin{align}
C^{(0)} = [0.5,\ 0.5]
\qquad \text{and} \qquad
C^0 = [0.5,\ 0.5].
\end{align}
They were equal in value, but they still played different roles:
\begin{align*}
C^{(0)} &: \text{the current cost at the start of training}, \\
C^0 &: \text{the fixed center of the global trust-region ball.}
\end{align*}

\subsection{Final Learned Cost Vector}
The final learned cost vector is the last accepted cost vector after the outer
BCD loop ends:
\begin{align}
C^{\star} = \texttt{final learned cost vector}.
\end{align}
In your reported run,
\begin{align}
C^\star \approx [0.7650,\ 0.9240].
\end{align}

\section{How the Cost Update Is Evaluated}
With $\alpha$ and $\beta$ fixed, the code computes the classwise loss vector
\begin{align}
L(\alpha,\beta)
=
\begin{bmatrix}
L_0(\alpha,\beta) \\
L_1(\alpha,\beta)
\end{bmatrix},
\end{align}
where
\begin{align}
L_k(\alpha,\beta)
=
\frac{1}{|pf|}
\sum_{m \in pf}
\frac{1}{n_m}
\sum_{j \in G_m : y_{m,j} = k}
\ell(z_{m,j}, y_{m,j}).
\end{align}

\paragraph{Important sign convention.}
Block A is a \emph{minimization} over $(\alpha,\beta)$ for fixed $C$. The cost
block is then treated through the reduced objective
\begin{align}
F(C) = \min_{\alpha,\beta} \Phi(\alpha,\beta;C),
\end{align}
and the trust-region cost step tries to find a direction that
\emph{increases} this reduced objective with respect to $C$. That is why the
BCD summary reports \texttt{Predicted increase} and \texttt{Actual increase}
rather than decreases.

The trust-region subproblem then proposes a step $d_t$ and a trial cost
\begin{align}
C_{\mathrm{trial}} = C^{(t)} + d_t.
\end{align}

\subsection{Predicted Increase}
Because the local model for the cost block is linear, the predicted increase is
\begin{align}
\mathrm{Pred}_t = L_t^\top d_t.
\end{align}
In the results table, this is the column called \texttt{Pred} or
\texttt{Pred inc}.

\subsection{Current Objective and Trial Objective}
Define the reduced objective
\begin{align}
F(C) = \min_{\alpha,\beta} \Phi(\alpha,\beta;C).
\end{align}

Then:
\begin{align*}
\text{current objective} &= F(C^{(t)}), \\
\text{trial objective} &= F(C_{\mathrm{trial}}).
\end{align*}

In code, both are computed approximately by solving Block A:
\begin{itemize}
\item the current objective is the latest accepted fixed-cost Block A value,
\item the trial objective is the fixed-cost Block A value after re-solving at
the trial cost.
\end{itemize}

\paragraph{What is the trust-region calculation in code?}
The code does not jump directly from $C^{(t)}$ to a closed-form optimal
$C^\star$. Instead it computes:
\begin{enumerate}
\item the class-loss vector $L_t = L(\alpha,\beta)$ with the current fixed
$\alpha$ and $\beta$,
\item the normalized direction
\begin{align}
\mathrm{direction}_t = \frac{L_t}{\|L_t\|_2},
\end{align}
\item a feasible step length $\tau_t$ limited by
\begin{align*}
\text{working trust-region limit}, \\
\text{global baseline-ball limit}, \\
\text{nonnegativity limit},
\end{align*}
\item the actual cost step
\begin{align}
d_t = \tau_t \frac{L_t}{\|L_t\|_2},
\end{align}
\item the trial cost
\begin{align}
C_{\mathrm{trial}} = C^{(t)} + d_t.
\end{align}
\end{enumerate}

So the trust-region step is a \emph{bounded move in cost space}. The final
learned cost is obtained only after repeating this outer BCD logic and
accepting or rejecting the trial cost moves.

\subsection{Actual Increase}
The actual increase is
\begin{align}
\mathrm{Ared}_t
=
F(C_{\mathrm{trial}}) - F(C^{(t)}).
\end{align}
In the results table, this is the column called \texttt{Actual} or
\texttt{Actual inc}.

\subsection{Ratio and Acceptance}
The trust-region ratio is
\begin{align}
\rho_t = \frac{\mathrm{Ared}_t}{\mathrm{Pred}_t}.
\end{align}

The code accepts the trial cost step when
\begin{align}
\rho_t \ge \eta_{\mathrm{acc}},
\end{align}
where
\begin{align}
\eta_{\mathrm{acc}} = \texttt{accept\_threshold}.
\end{align}

So:
\begin{align*}
\texttt{Accepted=True}
&\Longrightarrow
\text{move to } (C_{\mathrm{trial}}, \alpha_{\mathrm{trial}}, \beta_{\mathrm{trial}}), \\
\texttt{Accepted=False}
&\Longrightarrow
\text{reject the trial step and keep the previous } (C,\alpha,\beta).
\end{align*}

The summary field \texttt{Accepted cost steps} is simply the count of outer BCD
iterations for which \texttt{Accepted=True}.

\subsection{Why the Run Stopped After Four BCD Entries}
Although the run allowed up to $10$ outer iterations, it recorded only $4$
outer entries because the loop terminated early when no meaningful new cost step
was available. That is why the last line showed
\begin{verbatim}
Accepted: False | Termination: no_cost_step
\end{verbatim}
This means the training did \emph{not} use all ten outer iterations; it stopped
earlier because the cost-update step had essentially vanished.

\subsection{Run-Specific Cost Sequence}
Your run produced the following accepted cost path:
\begin{align}
C^{(0)} &= [0.5000,\ 0.5000], \\
C^{(1)} &= [0.6424,\ 0.6404], \\
C^{(2)} &= [0.7422,\ 0.9051], \\
C^{(3)} &= [0.7650,\ 0.9240].
\end{align}

Then the next outer iteration produced no meaningful cost step, so the training
terminated and kept
\begin{align}
C^\star = C^{(3)} = [0.7650,\ 0.9240].
\end{align}

This is why the BCD summary showed:
\begin{itemize}
\item $4$ recorded outer BCD entries,
\item $3$ accepted cost steps,
\item $1$ final termination with \texttt{no\_cost\_step}.
\end{itemize}

\paragraph{How many optimization rounds were done?}
For this run:
\begin{itemize}
\item the outer BCD loop recorded $4$ entries,
\item the initial fixed-cost Block A solve used $40$ Block A iterations,
\item each accepted trial cost step was followed by another fixed-cost Block A
solve of $40$ iterations,
\item the final outer iteration terminated before another trial Block A solve
was needed.
\end{itemize}

\paragraph{Which trust-region parameters produced this path?}
The run used:
\begin{align}
\Delta_1 &= \texttt{initial\_trust\_region\_radius} = 0.02, \\
\Delta_{\max} &= \texttt{max\_trust\_region\_radius} = 0.125, \\
\eta_{\mathrm{acc}} &= \texttt{accept\_threshold} = 0.25, \\
\eta_{\mathrm{expand}} &= \texttt{expand\_threshold} = 0.75, \\
\gamma_{\mathrm{shrink}} &= \texttt{shrink\_factor} = 0.25, \\
\gamma_{\mathrm{expand}} &= \texttt{expand\_factor} = 2.0.
\end{align}

So the final learned cost vector was not assumed in advance. It emerged from:
\begin{enumerate}
\item the starting cost $C^{(0)}$,
\item the fixed baseline center $C^0$,
\item the class-loss vectors $L_t$ computed after solving for $\alpha$ and
$\beta$,
\item the trust-region step $d_t$,
\item the accept/reject test based on the ratio $\rho_t$.
\end{enumerate}

\subsection{How to Read the BCD Summary Columns}
Each BCD row corresponds to \emph{one proposed outer cost update}. The row is
not reporting classification performance. It is reporting whether the
trust-region cost step for that outer iteration was good enough to keep.

Starting from the current accepted cost vector $C^{(t)}$, the code:
\begin{enumerate}
\item solves for the current $\alpha$ and $\beta$ at that cost,
\item computes the class-loss vector $L_t$,
\item proposes a bounded cost step $d_t$,
\item forms the trial cost
\begin{align}
C_{\mathrm{trial}} = C^{(t)} + d_t,
\end{align}
\item re-solves Block A at $C_{\mathrm{trial}}$,
\item compares what the local model predicted with what actually happened.
\end{enumerate}

The BCD summary fields mean:
\begin{align}
\text{current objective} &= F(C^{(t)}), \\
\text{trial objective} &= F(C_{\mathrm{trial}}), \\
\text{actual increase} &= F(C_{\mathrm{trial}}) - F(C^{(t)}), \\
\text{predicted increase} &= L_t^\top d_t, \\
\rho_t &= \frac{\text{actual increase}}{\text{predicted increase}}.
\end{align}

Here $F(C)$ is the reduced objective
\begin{align}
F(C) = \min_{\alpha,\beta} \Phi(\alpha,\beta;C),
\end{align}
so:
\begin{itemize}
\item \textbf{current objective} is the current fixed-cost Block A objective at
the accepted cost $C^{(t)}$,
\item \textbf{trial objective} is the fixed-cost Block A objective after
switching to the trial cost and re-solving $\alpha$ and $\beta$,
\item \textbf{predicted increase} is what the trust-region linear model expected
the cost move to achieve before re-solving Block A,
\item \textbf{actual increase} is what was actually achieved after re-solving
Block A at the trial cost,
\item \textbf{ratio} measures how accurate that prediction was.
\end{itemize}

\paragraph{How to interpret the ratio.}
\begin{align*}
\rho_t > 1 &:\ \text{the step worked better than predicted,} \\
\rho_t \approx 1 &:\ \text{the prediction was very accurate,} \\
0 < \rho_t < 1 &:\ \text{the step helped, but less than predicted,} \\
\rho_t < \eta_{\mathrm{acc}} &:\ \text{reject the trial cost step.}
\end{align*}

The \textbf{radius} field is the working trust-region radius. When the summary
prints
\begin{verbatim}
Radius: a -> b
\end{verbatim}
it means:
\begin{itemize}
\item $a$ was the radius used to construct the current trial step,
\item $b$ is the radius that will be used in the next outer iteration.
\end{itemize}

\paragraph{Run-specific interpretation of your BCD rows.}
\begin{itemize}
\item \textbf{Iteration 1}:
\begin{align*}
\text{current objective} &= 0.301791, \\
\text{trial objective} &= 0.433324, \\
\text{actual increase} &= 0.433324 - 0.301791 = 0.131533, \\
\text{predicted increase} &= 0.080115, \\
\rho_1 &= 1.641803.
\end{align*}
This means the first cost step worked \emph{better} than the local model
predicted. So the step was accepted, and the radius expanded from $0.02$ to
$0.04$.

\item \textbf{Iteration 2}:
\begin{align*}
\text{current objective} &= 0.433324, \\
\text{trial objective} &= 0.474506, \\
\text{actual increase} &= 0.041182, \\
\text{predicted increase} &= 0.140319, \\
\rho_2 &= 0.293487.
\end{align*}
This step still improved the reduced objective, but much less than predicted.
It was still accepted because the ratio remained above the acceptance threshold
$0.25$.

\item \textbf{Iteration 3}:
\begin{align*}
\text{current objective} &= 0.474506, \\
\text{trial objective} &= 0.485248, \\
\text{actual increase} &= 0.010741, \\
\text{predicted increase} &= 0.011537, \\
\rho_3 &= 0.931015.
\end{align*}
This means the local model was very accurate. The step was accepted, and the
radius expanded from $0.04$ to $0.08$.

\item \textbf{Iteration 4}:
\begin{align*}
\text{current objective} &= 0.485248, \\
\text{trial objective} &= 0.485248, \\
\text{actual increase} &= 0, \\
\text{predicted increase} &= 0.
\end{align*}
No meaningful new cost step was available, so the algorithm terminated with
\texttt{no\_cost\_step}. The current accepted cost vector was kept unchanged.
\end{itemize}

\paragraph{Important note.}
The objective in the BCD table is \emph{not} accuracy and it is \emph{not} just
plain BCE. It is the reduced fixed-cost training objective after optimizing
$\alpha$ and $\beta$, which includes the group cost-weighted BCE and the beta
regularization term.

\section{How the Radius Changes}
The working trust-region radius changes according to the ratio:
\begin{align}
\Delta_{t+1}
=
\begin{cases}
\gamma_{\mathrm{shrink}} \Delta_t, & \rho_t < \eta_{\mathrm{acc}}, \\
\min(\gamma_{\mathrm{expand}}\Delta_t,\Delta_{\max}), & \rho_t \ge \eta_{\mathrm{expand}}, \\
\Delta_t, & \text{otherwise.}
\end{cases}
\end{align}

So the \texttt{Radius: a -> b} field in the report means:
\begin{itemize}
\item $a$ is the working radius used to build the current trial cost step,
\item $b$ is the radius that will be used in the next outer BCD iteration.
\end{itemize}

\section{What Block A Reports Mean}
The Block A report lines such as
\begin{verbatim}
Initial Block A iter 0001 | Obj: ... | Group Cost-BCE: ... | Beta L1: ...
\end{verbatim}
describe the fixed-cost alpha/beta optimization. They do \emph{not} update the
cost vector.

\subsection{Group Cost-BCE}
\texttt{Group Cost-BCE} is the mean cost-weighted BCE across the overlapping
optimization groups:
\begin{align}
\frac{1}{|pf|}
\sum_{m \in pf}
f_m^C(\alpha,\beta,C).
\end{align}

\subsection{Sample Cost-BCE}
\texttt{Sample Cost-BCE} is the mean cost-weighted BCE computed on the full set
of valid sample-level predictions after stitching together the exact-profile
predictions.

\subsection{Beta L1}
\texttt{Beta L1} in the report is the weighted beta regularization term
\begin{align}
\lambda_\beta \sum_i \|\beta^i\|_1.
\end{align}

\subsection{Alpha Max L1}
\texttt{Alpha Max L1} is the largest $\ell_1$ norm among all profile-specific
alpha vectors:
\begin{align}
\max_m \|\alpha_m\|_1.
\end{align}

\section{What the Prediction Step Actually Does}
Training uses overlapping optimization groups $G_m$, but prediction does not.

\subsection{Training View}
During training, a sample can contribute to multiple optimization groups because
\begin{align}
G_m = \{j : (p_j \mathbin{\&} m) = m \}.
\end{align}
So the optimization-group tables are a \emph{training view} of the data.

\subsection{Prediction View}
At prediction time, each sample belongs to exactly one exact observed-source
profile:
\begin{align}
A_j = \{i : \text{source } i \text{ is present for sample } j\}.
\end{align}
The prediction for sample $j$ is then computed only from that exact profile:
\begin{align}
z_j
=
\sum_{i \in A_j}
\alpha_{A_j}^i x_j^i \beta^i,
\qquad
\hat{p}_j = \sigma(z_j).
\end{align}
The class prediction is
\begin{align}
\hat{y}_j =
\begin{cases}
1, & \hat{p}_j \ge \tau, \\
0, & \hat{p}_j < \tau,
\end{cases}
\end{align}
where
\begin{align}
\tau = \texttt{decision\_threshold}.
\end{align}
In your run, $\tau = 0.5$.

\paragraph{Meaning of the decision threshold.}
The model first outputs a probability $\hat{p}_j$. The decision threshold is
the cutoff that turns that probability into a hard class label. It does not
change how the model is trained; it changes how probabilities are converted into
\texttt{0}/\texttt{1} predictions for the metrics table.

\section{How to Read the Alpha Table}
The learned alpha table is indexed by source profiles such as
\begin{align*}
(\texttt{source1}),\ 
(\texttt{source1},\texttt{source2}),\ 
(\texttt{source1},\texttt{source2},\texttt{source3}).
\end{align*}

Each row means:
\begin{quote}
\emph{If a sample has exactly this source-availability pattern at prediction
time, how should the model weight the participating sources?}
\end{quote}

For example, the row
\begin{verbatim}
('source1','source3'): {'source1': 0.075..., 'source3': 0.924...}
\end{verbatim}
means that, for the profile where sources $1$ and $3$ are present and source
$2$ is absent, the predictor relies far more on source $3$ than source $1$.

\section{How to Read the Beta Table}
There is one beta vector per source:
\begin{align*}
\beta^{(1)} &\in \mathbb{R}^{10}, \\
\beta^{(2)} &\in \mathbb{R}^{8}, \\
\beta^{(3)} &\in \mathbb{R}^{25}.
\end{align*}
These lengths match the feature counts in the three source CSV files.

The beta summary table reports:
\begin{itemize}
\item \textbf{shape}: the number of coefficients in that source,
\item \textbf{$\ell_2$ norm}: overall coefficient magnitude,
\item \textbf{$\ell_1$ norm}: total absolute magnitude, directly tied to the
regularizer,
\item \textbf{nonzero coefficients}: how many source features remain active
after the proximal $\ell_1$ shrinkage.
\end{itemize}

\section{How to Read the Two Profile-Based Result Tables}
The report prints two different profile tables because they answer two different
questions.

\subsection{Exact-Profile Training Metrics}
These rows are \emph{disjoint}. Each sample belongs to exactly one exact profile
such as
\begin{align*}
001,\ 010,\ 011,\ 100,\ 101,\ 110,\ 111.
\end{align*}
So this table is the closest to the prediction-time view of the model.

\subsection{Optimization-Group Training Metrics}
These rows are the overlapping training groups used by iSFS during fitting. A
sample may appear in more than one of these groups. So this table is the
closest to the training-time decomposition of the objective.

\paragraph{Simple interpretation.}
\begin{align*}
\text{Exact-profile table} &:\ \text{``How does the final model perform on each exact missingness pattern?''} \\
\text{Optimization-group table} &:\ \text{``How does the final model behave on each overlapping iSFS training group?''}
\end{align*}

\section{Metric Definitions That Appear in the Report}
\begin{longtable}{>{\raggedright\arraybackslash}p{0.22\textwidth} >{\raggedright\arraybackslash}p{0.68\textwidth}}
\toprule
\textbf{Metric} & \textbf{Meaning} \\
\midrule
\endhead
Accuracy & Fraction of all predictions that are correct. \\
Precision & Among predicted positives, the fraction that are truly positive. \\
Recall / Sensitivity & Among true positives, the fraction correctly detected. \\
Specificity & Among true negatives, the fraction correctly detected. \\
F1 score & Harmonic mean of precision and recall. \\
Balanced accuracy & Average of recall and specificity. \\
MCC & Matthews correlation coefficient. This uses TP, TN, FP, and FN together and is often more informative than accuracy when the classes are not perfectly balanced. \\
Brier score & Mean squared error of the predicted probabilities:
$\frac{1}{n}\sum_j (\hat{p}_j - y_j)^2$. Lower is better. It measures probability calibration and probability accuracy. \\
Log loss & Binary cross-entropy on the final predicted probabilities. Lower is better. \\
Probability range & Minimum and maximum predicted probability over the evaluated set. This tells you how conservative or extreme the model's probability outputs were. \\
\bottomrule
\end{longtable}

\section{Where These Objects Live in the Code}
\begin{longtable}{>{\raggedright\arraybackslash}p{0.28\textwidth} >{\raggedright\arraybackslash}p{0.62\textwidth}}
\toprule
\textbf{Object} & \textbf{Main implementation location} \\
\midrule
\endhead
Real-data loading & \texttt{run\_training.py}, function \texttt{load\_real\_dataset} \\
Outer BCD loop & \texttt{train\_model.py}, function \texttt{train\_blockwise\_logistic\_model} \\
Fixed-cost Block A loop & \texttt{train\_model.py}, function \texttt{solve\_block\_a\_with\_fixed\_cost} \\
Alpha subproblem solver & \texttt{alpha\_update.py}, function \texttt{solve\_alphas\_with\_fixed\_betas} \\
Beta subproblem solver & \texttt{beta\_update.py}, function \texttt{solve\_betas\_with\_fixed\_alphas} \\
Cost trust-region step & \texttt{costs.py}, function \texttt{solve\_cost\_trust\_region\_subproblem} \\
Acceptance ratio and radius update & \texttt{costs.py}, functions \texttt{compute\_trust\_region\_ratio} and \texttt{update\_trust\_region\_radius} \\
Exact profiles & \texttt{missingness\_profiles.py}, function \texttt{build\_exact\_profiles} \\
Overlapping optimization groups & \texttt{missingness\_profiles.py}, function \texttt{build\_missingness\_profiles} \\
Prediction logits and probabilities & \texttt{prediction.py}, functions \texttt{compute\_all\_logits}, \texttt{predict\_proba}, and \texttt{predict\_class} \\
Metrics in the final report & \texttt{run\_training.py}, functions \texttt{compute\_binary\_metrics}, \texttt{compute\_training\_summary}, and \texttt{compute\_profile\_metrics} \\
\bottomrule
\end{longtable}

\section{Final Summary}
The key point is that the run report mixes together:
\begin{enumerate}
\item \textbf{model quantities} such as $\alpha$, $\beta$, $C$, and the
profile groups,
\item \textbf{solver controls} such as learning rates and maximum iteration
counts,
\item \textbf{trust-region diagnostics} such as predicted increase, actual
increase, ratio, and radius,
\item \textbf{evaluation outputs} such as accuracy, MCC, Brier score, and the
exact-profile / optimization-group tables.
\end{enumerate}

Those are all part of the same training script, but they belong to different
layers of the implementation. The outer iteration counts, Block A iteration
counts, and alpha/beta subproblem iteration counts are \emph{solver controls}.
The cost vectors, alpha weights, and beta coefficients are \emph{learned model
objects}. The trust-region ratio and radius are \emph{cost-update diagnostics}.
The metrics tables are \emph{evaluation summaries} of the final fitted model.

\end{document}
